\section{实值函数族的包络}

记号约定:$E$是集合,$\left(f_{i}\right)_{i\in I}$是$\Hom_{\mathsf{Set}}\left(E,\overline{\R}\right)$的元素族.

\begin{definition}{实值函数族的上包络和下包络}{}
	\begin{enumerate}
		\item $\left(f_{i}\right)_{i\in I}$在$\Hom_{\mathsf{Set}}\left(E,\overline{\R}\right)$中的上确界记作$\sup\limits_{i\in I}f_{i}$,称为$\left(f_{i}\right)_{i\in I}$的上包络.
		\item $\left(f_{i}\right)_{i\in I}$在$\Hom_{\mathsf{Set}}\left(E,\overline{\R}\right)$中的下确界记作$\inf\limits_{i\in I}f_{i}$,称为$\left(f_{i}\right)_{i\in I}$的下包络.
	\end{enumerate}
\end{definition}

\begin{proposition}{}{}
	记$\mathfrak{F}\left(E\right)=\fin\left(E\right)\setminus\left\{\emptyset\right\}$.给$\Hom_{\mathsf{Set}}\left(E,\overline{\R}\right)$装备乘积拓扑,那么$\sup\limits_{i\in I}f_{i}=\limit_{H\in\mathfrak{F}\left(E\right)}\sup\limits_{i\in H}f_{i}$.
\end{proposition}

\begin{definition}{一致上有界,一致下有界,一致有界}{}
	\begin{enumerate}
		\item 如果$\exists a\in\R\forall x\in E\forall i\in I\left(f_{i}\left(x\right)\leq a\right)$,则称$\left(f_{i}\right)_{i\in I}$在$E$上一致上有界.
		\item 如果$\exists a\in\R\forall i\in I\forall x\in E\left(f_{i}\left(x\right)\geq a\right)$,则称$\left(f_{i}\right)_{i\in I}$在$E$上一致下有界.
		\item 如果$\left(f_{i}\right)_{i\in I}$在$E$上既一致上有界,又一致下有界,则称$\left(f_{i}\right)_{i\in I}$在$E$上一致有界.
	\end{enumerate}
\end{definition}

\begin{remark}{}{}
	\begin{enumerate}
		\item $\left(f_{i}\right)_{i\in I}$在$E$上一致上有界$\iff$ $\sup\limits_{i\in I}f_{i}$在$E$上有上界.
		\item $\left(f_{i}\right)_{i\in I}$在$E$上一致有界$\iff$ $\left(\left|f_{i}\right|\right)_{i\in I}$在$E$上有上界.
	\end{enumerate}
\end{remark}